\section{Implementation and verification details}
\label{app:numerics}

\subsection{Error budget at m=4}

The independently recomputed Bayes-risk enclosure width is
$4.9777\times10^{-11}$.  Of this, the reported Simpson-remainder width is
$4.9696\times10^{-11}$, the Gaussian tail upper bound is
$8.084\times10^{-14}$, and arithmetic rounding in the Simpson sum is below
$3.1\times10^{-40}$.  At risk points encountered by separation, the largest
fourth-order remainder width is $6.5690\times10^{-11}$ and the largest tail
upper bound is $8.084\times10^{-14}$.  The branch envelope width is
$5.0\times10^{-9}$.

\subsection{Independent reconstruction}

The discovery solver computes posterior states with vectorized NumPy and a
specialized flow tube.  The final verifier instead:
\begin{enumerate}
\item constructs Arb balls from decimal strings;
\item evaluates each mixture exponential directly;
\item forms posterior central moments through order five;
\item implements a separate Simpson panel loop;
\item uses direct interval boxes for curvature rather than the common-mean
flow formula; and
\item rebuilds its own branch tree from $[0,m]$.
\end{enumerate}
The two paths share the mathematical definition of risk while using different
numerical representations and control flow.

\subsection{Failure modes examined}

Independent intervals for $f_0$, $f_1^2$, and $1/f_0$ repeatedly count the
same uncertainty and become particularly wide in the tails.  A global
Lipschitz bound scales only linearly with cell width and exhausted the
5001-node budget even at tolerance $10^{-4}$.  A global curvature bound also
exhausted that budget.  One-noise affine arithmetic helps on narrow cells but
was sometimes wider than the posterior-flow representation.  These empirical
results characterize the tested representations and tolerances.
